%==============================================================================
% DevPrism Beamer presentation scaffold
% Compile locally with:   tectonic slides.tex
% (or, with a TeXLive install:  pdflatex slides.tex   run twice for the TOC)
%
% Uses only the built-in "Madrid" theme + standard TeXLive packages.
% Optional metropolis theme is commented out and guarded below.
%==============================================================================
\documentclass[aspectratio=169]{beamer}  % 16:9 widescreen; use 43 for 4:3

%------------------------------------------------------------------------------
% Theme
%------------------------------------------------------------------------------
% Default: clean built-in theme, no extra fonts required.
\usetheme{Madrid}
\usecolortheme{default}
\setbeamertemplate{navigation symbols}{}  % hide the navigation clutter

% --- Optional modern theme (requires the "metropolis" package) --------------
% To use it instead of Madrid: comment out the three lines above and
% uncomment the block below. metropolis is not part of the minimal TeXLive
% scheme, so leave Madrid unless you know it is installed.
%
% \usetheme{metropolis}
% \metroset{block=fill}

%------------------------------------------------------------------------------
% Packages (all standard TeXLive; guarded so missing ones do not abort)
%------------------------------------------------------------------------------
\usepackage[utf8]{inputenc}
\usepackage[T1]{fontenc}
\usepackage{graphicx}   % \includegraphics
\usepackage{booktabs}   % nice tables: \toprule \midrule \bottomrule

% TikZ is optional: only used for the self-contained figure placeholder.
% It ships with standard TeXLive. If it is missing, the figure frame falls
% back to a plain framebox (see \DrawPlaceholder below).
\newif\ifhastikz
\IfFileExists{tikz.sty}{\usepackage{tikz}\hastikztrue}{\hastikzfalse}

% Define the figure placeholder ONCE in the preamble, where \if..\fi is safe.
% (Raw \if..\else..\fi inside a beamer frame body can confuse overlay parsing,
% so we hide it behind a command.) Replace the call in the figure frame with
% your own \includegraphics when you have a real image.
\ifhastikz
  \newcommand{\figureplaceholder}{%
    \begin{tikzpicture}
      \draw[thick, rounded corners] (0,0) rectangle (7,4);
      \node at (3.5,2.2) {\large Replace with \texttt{\textbackslash includegraphics}};
      \node at (3.5,1.4) {\small put \texttt{figure.png} next to this file};
    \end{tikzpicture}}
\else
  \newcommand{\figureplaceholder}{%
    \framebox[0.7\textwidth][c]{\rule{0pt}{4cm}Figure placeholder}}
\fi

%------------------------------------------------------------------------------
% Title metadata  -- EDIT THESE
%------------------------------------------------------------------------------
\title[Short Title]{Your Talk Title}
\subtitle{An optional subtitle}
\author[A.~Author]{Ada Author}
\institute[Inst.]{Your Lab / Institution}
\date{\today}

%------------------------------------------------------------------------------
% Auto table of contents at the start of every section
%------------------------------------------------------------------------------
\AtBeginSection[]{%
  \begin{frame}{Outline}
    \tableofcontents[currentsection]
  \end{frame}%
}

\begin{document}

%==============================================================================
% Title frame
%==============================================================================
\begin{frame}[plain]
  \titlepage
\end{frame}

%==============================================================================
% Outline frame
%==============================================================================
\begin{frame}{Outline}
  \tableofcontents
\end{frame}

%==============================================================================
\section{Motivation}
%==============================================================================
\begin{frame}{Why this matters}
  \begin{itemize}
    \item One short phrase per bullet.
    \item State the problem, not the whole paper.
    \item End with the question your talk answers.
  \end{itemize}
\end{frame}

%==============================================================================
\section{Method}
%==============================================================================

%------------------------------------------------------------------------------
% Two-column frame: text on the left, list/figure slot on the right
%------------------------------------------------------------------------------
\begin{frame}{Approach (two columns)}
  \begin{columns}[T]  % [T] = top-align the columns
    \begin{column}{0.48\textwidth}
      \textbf{Idea}
      \begin{itemize}
        \item Describe the method in words.
        \item Keep it to a few phrases.
        \item Point to the figure on the right.
      \end{itemize}
    \end{column}
    \begin{column}{0.48\textwidth}
      \textbf{Pipeline}
      \begin{enumerate}
        \item Input data
        \item Transform / model
        \item Output / result
      \end{enumerate}
    \end{column}
  \end{columns}
\end{frame}

%------------------------------------------------------------------------------
% Figure frame: replace the placeholder with \includegraphics{your-figure}
%------------------------------------------------------------------------------
\begin{frame}{A key figure}
  \begin{center}
    % --- Real usage (uncomment and set your filename): -------------------
    % \includegraphics[width=0.7\textwidth]{figure}
    %
    % --- Self-contained placeholder so the template compiles with no image:
    \figureplaceholder
  \end{center}
  \begin{center}\small Takeaway: state the one thing this figure shows.\end{center}
\end{frame}

%==============================================================================
\section{Results}
%==============================================================================

%------------------------------------------------------------------------------
% Bulleted list with \pause overlays (step reveal) + \only / \onslide / \alert
%------------------------------------------------------------------------------
\begin{frame}{Main results (revealed step by step)}
  \begin{itemize}
    \item First finding appears immediately.
    \pause
    \item Second finding appears on the next click.
    \pause
    \item \alert<4>{Third finding is highlighted on click 4.}
  \end{itemize}

  \vspace{0.5em}
  % \only<n-> shows ONLY on step n and later (takes no space before).
  \only<5->{\textbf{Summary line that appears last.}}

  % \onslide<n-> reserves the space even before it appears (no jumping).
  \onslide<6->{\par\small (Reserved-space note shown on the final step.)}
\end{frame}

%------------------------------------------------------------------------------
% A compact results table (lift a 3-4 row summary from the paper)
%------------------------------------------------------------------------------
\begin{frame}{Numbers at a glance}
  \begin{center}
    \begin{tabular}{lcc}
      \toprule
      Method        & Metric A & Metric B \\
      \midrule
      Baseline      & 0.71     & 0.65 \\
      \textbf{Ours} & \textbf{0.88} & \textbf{0.83} \\
      \bottomrule
    \end{tabular}
  \end{center}
  \begin{center}\small One sentence: what the table proves.\end{center}
\end{frame}

%==============================================================================
\section{Conclusion}
%==============================================================================
\begin{frame}{Key takeaways}
  \begin{itemize}
    \item Restate the answer to the opening question.
    \item Name the single most important result.
    \item State what comes next / future work.
  \end{itemize}
\end{frame}

%==============================================================================
% Closing / thanks frame
%==============================================================================
\begin{frame}[plain]
  \begin{center}
    {\Huge Thank you!}\\[1em]
    {\large Questions?}\\[2em]
    \small Ada Author \quad\textbar\quad ada@example.org
  \end{center}
\end{frame}

\end{document}
